\chapter{Literature Review}
\label{ch:litreview}

\section{Scope and method of the review}

This review covers the three bodies of work the project draws on: modelling of
photovoltaic devices, maximum power point tracking algorithms, and the control
of the power converter that carries out the tracking. It is organised
thematically rather than chronologically, because the interesting structure in
this literature is not a progression but a persistent \emph{separation}: the
device modelling and the algorithm literature developed largely independently
of the power-electronics control literature, and the gap between them is where
this project sits.

Each section states what is settled, what is contested, and what is assumed
without justification. Section~\ref{sec:gap} draws the threads together into the
gap this project addresses.

\section{Modelling the photovoltaic device}

\subsection{The single-diode model is settled; its parameterisation is not}

There is no serious disagreement that the single-diode equation with series and
shunt resistance is the right lumped model for a crystalline silicon module at
the terminal level. Two-diode models exist and improve accuracy at very low
irradiance by separating diffusion from recombination current, at the cost of
two additional parameters that datasheets do not constrain. For terminal-level
control studies the single-diode form is universal.

What is genuinely contested is how to obtain the five parameters from the four
or five numbers a datasheet provides. Three families of approach appear.

\begin{description}[leftmargin=0pt, style=unboxed]
\item[Fix the ideality factor and iterate.] Villalva, Gazoli and
  Filho~\cite{villalva} fix $a$ at an assumed value and iterate on $R_s$,
  computing $R_{sh}$ at each step from the requirement that the modelled peak
  power equal the datasheet peak power. The method is robust, easy to implement
  and has been cited thousands of times. Its weakness is that one parameter is
  chosen by the modeller rather than identified, and the temperature coefficient
  $\beta_{Voc}$, which the datasheet does provide, is never used.
\item[Solve a determined system.] De Soto, Klein and Beckman~\cite{desoto} close
  the system by requiring the model to reproduce the measured $\beta_{Voc}$,
  giving five equations in five unknowns. This uses all the available datasheet
  information and leaves nothing to judgement. It is the basis of the NREL
  System Advisor Model. Its weakness is numerical: the system is badly scaled,
  since $I_0$ spans ten orders of magnitude, and naive solvers fail to converge.
\item[Fit to measured curves.] Where a full measured I--V curve is available,
  least-squares or evolutionary fitting outperforms both. This is not available
  from a datasheet and so is not applicable to a design study.
\end{description}

\subsection{What the literature rarely checks}

Almost every published extraction reports agreement at $\Isc$, $\Voc$ and the
maximum power point, because those are the constraints imposed. Very few report
whether the \emph{slope} of the curve is right anywhere other than at those
points. For an energy-yield study that omission is harmless. For a control
study it is not, because as Chapter~\ref{ch:converter} shows, the small-signal
plant depends on $\mathrm{d}I/\mathrm{d}V$ and not on the datasheet points at
all. This project therefore adopts the De Soto closure, solves the badly scaled
system properly, and verifies the derivative independently.

\section{Maximum power point tracking algorithms}

\subsection{The canonical survey and its framing}

Esram and Chapman~\cite{esram} remains the standard survey, comparing nineteen
methods on cost, sensors required, convergence speed and whether they find a
true or approximate maximum. It established the vocabulary the field still uses.
Its framing, however, is algorithmic: methods are compared as search procedures,
with the converter treated as a means of imposing whatever operating point the
algorithm requests. The possibility that the converter dynamics might limit
which algorithms can work, or dictate their tuning, is largely outside its
scope.

\subsection{Perturb and observe and its known pathologies}

Perturb and observe dominates practice because it needs only two sensors and a
few lines of code. Its two well-documented weaknesses are the steady-state
limit cycle, whose amplitude is set by the perturbation size, and the drift or
confusion phenomenon under changing irradiance.

Femia, Petrone, Spagnuolo and Vitelli~\cite{femia} gave the first systematic
treatment of the second problem. Their central insight is directly
control-theoretic: the perturbation period must be long enough for the converter
to settle, otherwise the algorithm measures a transient rather than a new
steady state, and the sign of the observed power change becomes unreliable. They
derive a lower bound on the perturbation period from the converter's settling
behaviour. This paper is the closest antecedent to the framing of the present
project, and its influence on Chapter~\ref{ch:control} is explicit.

Their later work~\cite{femia_book} extends this into a book-length treatment
that consistently treats the converter as part of the tracking problem rather
than as a black box.

\subsection{Incremental conductance: a claimed advantage that is hard to observe}

Incremental conductance~\cite{hussein} was introduced specifically to address
rapidly changing atmospheric conditions, on the argument that evaluating the
optimality condition $\mathrm{d}I/\mathrm{d}V = -I/V$ directly is more reliable
than differencing power. The claim of superiority is widely repeated.

The most careful comparison is by Sera and colleagues~\cite{sera}, who conclude
that under equivalent implementation the two methods are essentially equivalent,
because both ultimately compute the sign of the same quantity, and that reported
differences usually come from unequal step sizes or unequal sampling rather than
from the algorithms themselves. Elgendy, Zahawi and Atkinson~\cite{elgendy}
reach a similar conclusion experimentally, finding implementation details
dominate the choice of method.

The present project reproduces exactly this result, and Chapter~\ref{ch:mppt}
explains why: with the same step and the same sample period the two algorithms
compute $\mathrm{sign}(\mathrm{d}P/\mathrm{d}V)$ by two routes. The distinction
should reappear when measurement noise is modelled, which this project does not
do and which is stated as a limitation.

\subsection{Variable-step methods and their risk}

A large literature proposes making the perturbation size adaptive, typically
proportional to $|\mathrm{d}P/\mathrm{d}V|$, so that the tracker moves quickly
when far from the peak and finely when near it. These papers almost invariably
demonstrate improvement on a step-change benchmark.

What is much less often examined is the stability of the adaptation law itself.
The adaptation is a feedback loop, and Chapter~\ref{ch:benchmark} of this report
shows it can be a positively fed-back one: if drift sends the tracker the wrong
way, $|\mathrm{d}P/\mathrm{d}V|$ grows, which enlarges the step, which
accelerates the departure. The result observed here, that the best performer on
a step test is the worst on a ramp test, suggests that step-change benchmarking
is actively misleading for this class of method.

\subsection{Extremum seeking and ripple correlation control}

Extremum seeking control has a rigorous foundation independent of photovoltaics.
Krsti\'{c} and Wang~\cite{krstic} established local stability for general
nonlinear plants using averaging and singular perturbation, and Ariyur and
Krsti\'{c}~\cite{ariyur} give the design treatment. The method is attractive
here precisely because it does not require knowing the location of the optimum,
only that one exists.

Leyva and colleagues~\cite{leyva} applied it to photovoltaic tracking and
established the design constraints: the dither must be slow enough for the
converter's voltage loop to follow it and fast enough to be separable from
genuine irradiance variation. The permanent dither is an unavoidable
steady-state cost.

Ripple correlation control~\cite{rcc} is the elegant refinement: instead of
adding a dither, it uses the switching ripple the converter already produces,
correlating the ripple in power with the ripple in voltage to estimate the
gradient. The steady-state penalty disappears because the excitation was there
anyway. It requires bandwidth sufficient to observe the ripple, which is why it
demands a switched rather than an averaged model, and it is identified in
Chapter~\ref{ch:conclusions} as the natural next step for this project.

\subsection{Global methods for partial shading}

Under non-uniform illumination the P--V characteristic acquires multiple local
maxima, and any local hill-climber can converge to the wrong one. Nguyen and
Low~\cite{nguyen} are representative of the global-search response. This
literature is important but orthogonal to the present work, which assumes
uniform illumination throughout and says so.

\section{Control of the tracking converter}

\subsection{Averaging is the standard tool}

State-space averaging, due to Middlebrook and \'{C}uk~\cite{middlebrook} and
developed systematically by Erickson and Maksimovi\'{c}~\cite{erickson},
replaces the switched converter with a continuous model valid well below the
switching frequency. It is the foundation of essentially all converter control
design and is used unmodified here. What matters for this project is a detail
that general treatments do not emphasise: when the source is a photovoltaic
array rather than a voltage source, the source's dynamic resistance enters the
averaged model and varies over orders of magnitude with operating point.

\subsection{The right half plane zero, and where it actually is}

The boost converter's non-minimum-phase duty-to-output-voltage response is
textbook material~\cite{erickson} and is the standard explanation for why boost
converters have limited bandwidth. In the photovoltaic literature this is
sometimes cited as a general limitation on MPPT bandwidth.

That transfer is not, however, the one a photovoltaic-side tracker uses.
Chapter~\ref{ch:converter} derives the duty-to-\emph{input}-voltage response and
shows it has no finite zero at all. The apparent limitation disappears once the
correct controlled variable is identified. This point is made in the
power-electronics literature but is frequently lost in MPPT papers.

\subsection{Cascade control of photovoltaic converters}

Current-mode and cascaded control of DC--DC converters is long-established
practice. Bratcu and colleagues~\cite{bratcu} treat power optimisation in
cascaded photovoltaic converter systems explicitly, and Kjaer, Pedersen and
Blaabjerg~\cite{kjaer} survey the converter architectures in which such loops
sit. The benefits, damping of the output filter resonance and natural current
limiting, are well understood in the power-electronics community.

The gap is that MPPT papers rarely quantify them. It is common to see a
photovoltaic voltage loop described as ``a PI controller'' with gains stated but
no crossover frequency, no phase margin, and no indication of whether those
margins hold at other operating points. The present project computes those
margins across the whole envelope and compares the two architectures on that
basis.

\subsection{Benchmarking}

Comparisons in the MPPT literature use widely varying test signals: step
changes, ramps of unstated slope, or recorded irradiance traces. This makes
published efficiency figures difficult to compare. EN 50530~\cite{en50530}
exists precisely to fix this, specifying trapezoidal irradiance profiles with
defined slopes and dwell times and defining dynamic MPPT efficiency as the
figure of merit. Its adoption in the academic literature is patchy. This project
uses profiles of that form, and states its deviations from the standard
explicitly.

\section{Synthesis and the gap addressed}
\label{sec:gap}

Three observations follow from the review.

\begin{enumerate}[leftmargin=1.6em]
  \item \textbf{Device modelling and control analysis are usually decoupled.}
        Parameter extraction is validated against datasheet points, which
        constrain position but not curvature, while the small-signal plant
        depends only on curvature. A model can therefore be validated by the
        usual criteria and still be unfit for control design.
  \item \textbf{MPPT algorithms are compared as search procedures, not as
        outer loops of a controlled plant.} Sample periods, step sizes and
        stability limits are usually chosen empirically, when they are in fact
        determined by measurable properties of the converter loop.
  \item \textbf{Benchmarking practice under-tests the failure modes that
        matter.} Step-change tests reward exactly the adaptive schemes that
        ramp tests expose.
\end{enumerate}

The gap this project addresses is the connective tissue. Specifically, it aims
to:

\begin{itemize}[leftmargin=1.4em]
  \item identify a photovoltaic model whose \emph{derivatives} are verified, not
        only its operating points;
  \item derive the small-signal plant explicitly, quantify how it varies over
        the whole operating envelope, and show what that variation does to a
        conventional design;
  \item derive the MPPT sample period from the measured closed-loop settling
        time rather than choosing it;
  \item compare algorithms under identical inner loops on a standard-derived
        dynamic benchmark, so that differences are attributable to the
        algorithms;
  \item and validate the whole chain against a switched circuit model that makes
        none of the averaging assumptions.
\end{itemize}

None of the individual ingredients is novel. What the review suggests is
under-served is their combination into a single verified chain from datasheet to
switched circuit, with each link checked against something independent. That is
what Chapter~\ref{ch:method} sets out as the methodology.
